\begin{abstract}
Orbital missions have mapped the Moon with many instruments, from multispectral images taken under changing sunlight to static maps of topography, roughness and composition. The products are co-registered but differ in resolution, units and physics, and labels are scarce. We propose \method, a self-supervised model that learns from this stack by predicting, in representation space, what one instrument would record at a 60\,km footprint given another instrument's view of it. Per-product stems give every product the same ground cell per token, from 60\,m to 1\,km per pixel. A cross-attention predictor, told the target product and the sun under which it was observed, predicts the token grid of a moving-average encoder. We hypothesize that this task forces the static maps to encode what shapes the appearance of the surface, and turns the predictor into a sun-conditioned renderer in latent space. We also argue that the design has a failure mode the loss hides: a position code, in which tokens depend on their place in the grid rather than on the terrain. We show why the loss, within-image rank and VICReg penalties should not see it, and why a covariance penalty and a variance hinge on the predictions should push toward it. We then propose tile-centered cross-sample monitors that the position, product and location channels cannot inflate. Finally, we lay out experiments that test each hypothesis, from controlled pretraining ablations with probe-based ground truth to crater detection, shape and albedo from shading, and a sun-swap test of the predictor. \TODO{headline results}
\end{abstract}
